\documentclass[11pt,a4paper]{article}

% ── Codificación y fuentes ────────────────────────────────────────────────────
\usepackage[utf8]{inputenc}
\usepackage[T1]{fontenc}
\usepackage{lmodern}
\usepackage[spanish]{babel}

% ── Geometría y espaciado ─────────────────────────────────────────────────────
\usepackage[top=2.5cm, bottom=2.5cm, left=2.8cm, right=2.8cm]{geometry}
\usepackage{setspace}
\setstretch{1.15}
\usepackage{parskip}

% ── Tablas ────────────────────────────────────────────────────────────────────
\usepackage{booktabs}
\usepackage{tabularx}
\usepackage{array}
\usepackage{longtable}
\usepackage{enumitem}

% ── Colores y cajas ───────────────────────────────────────────────────────────
\usepackage[dvipsnames]{xcolor}
\usepackage{tcolorbox}
\tcbuselibrary{skins, breakable}

% ── Cabeceras y pies de página ────────────────────────────────────────────────
\usepackage{fancyhdr}
\usepackage{lastpage}
\pagestyle{fancy}
\fancyhf{}
\fancyhead[L]{\small\textbf{Informe de Evaluación} --- Física para Bioanalistas}
\fancyhead[R]{\small rosmeilys sucre 33386357}
\fancyfoot[C]{\small Página \thepage\ de \pageref{LastPage}}
\renewcommand{\headrulewidth}{0.4pt}

% ── Hiperenlaces ──────────────────────────────────────────────────────────────
\usepackage[colorlinks=true, linkcolor=NavyBlue, urlcolor=NavyBlue]{hyperref}

% ── Paleta de colores personalizados ─────────────────────────────────────────
\definecolor{desempeno_excelente}{HTML}{1A6B2F}
\definecolor{desempeno_bueno}{HTML}{2E5A9C}
\definecolor{desempeno_suficiente}{HTML}{8B6914}
\definecolor{desempeno_deficiente}{HTML}{C0392B}
\definecolor{desempeno_insuficiente}{HTML}{7B241C}
\definecolor{grisFondo}{HTML}{F5F5F5}
\definecolor{azulTitulo}{HTML}{1B2A6B}
\definecolor{verdeFortaleza}{HTML}{1A6B2F}
\definecolor{naranjaMejora}{HTML}{A04000}

% ── Estilos de cajas ─────────────────────────────────────────────────────────
\tcbset{
  cajaTitulo/.style={
    enhanced, breakable,
    colback=azulTitulo!8, colframe=azulTitulo,
    fonttitle=\bfseries\large, coltitle=white,
    attach boxed title to top left={yshift=-2mm, xshift=4mm},
    boxed title style={colback=azulTitulo},
    top=6pt, bottom=6pt, left=8pt, right=8pt,
  },
  cajaFortaleza/.style={
    enhanced, breakable,
    colback=verdeFortaleza!6, colframe=verdeFortaleza!60,
    fonttitle=\bfseries, coltitle=verdeFortaleza,
    top=4pt, bottom=4pt, left=8pt, right=8pt,
  },
  cajaMejora/.style={
    enhanced, breakable,
    colback=naranjaMejora!6, colframe=naranjaMejora!60,
    fonttitle=\bfseries, coltitle=naranjaMejora,
    top=4pt, bottom=4pt, left=8pt, right=8pt,
  },
  cajaRecomendacion/.style={
    enhanced, breakable,
    colback=grisFondo, colframe=gray!50,
    fonttitle=\bfseries, coltitle=black,
    top=4pt, bottom=4pt, left=8pt, right=8pt,
  },
}

% ─────────────────────────────────────────────────────────────────────────────
\begin{document}

% ══ PORTADA ══════════════════════════════════════════════════════════════════
\begin{center}
  {\Large\bfseries\color{azulTitulo} Universidad de Oriente/Núcleo Sucre --- Física para Bioanalistas}\\[4pt]
  {\large Informe de Evaluación Preliminar}\\[12pt]
  \rule{\linewidth}{1.2pt}\\[6pt]
  {\LARGE\bfseries rosmeilys sucre 33386357}\\[4pt]
  \rule{\linewidth}{0.6pt}
\end{center}

% ══ FICHA TÉCNICA ════════════════════════════════════════════════════════════
\vspace{4pt}
\begin{tcolorbox}[cajaTitulo, title=Datos del informe]
\begin{tabular}{@{}llll@{}}
  \textbf{Estudiante:}  & rosmeilys sucre 33386357  &
  \textbf{Fecha:}       & 2026-06-26 \\[4pt]
  \textbf{Archivo PDF:} & \texttt{rosmeilys\_sucre\_33386357.pdf}  &
  \textbf{Páginas:}     & 3 \\
\end{tabular}
\end{tcolorbox}

% ══ RESULTADO GLOBAL ═════════════════════════════════════════════════════════
\vspace{6pt}
\begin{center}
  \begin{tcolorbox}[
    enhanced,
    colback=white, colframe=desempeno_bueno,
    width=0.62\linewidth,
    boxrule=2pt,
    halign=center, valign=center,
    top=8pt, bottom=8pt,
  ]
    \centering
    {\large\bfseries Puntaje sugerido}\\[6pt]
    {\Huge\bfseries\color{desempeno_bueno} 8.3/10}\\[6pt]
    {\large Nivel de desempeño: \textbf{\color{desempeno_bueno} Bueno}}
  \end{tcolorbox}
\end{center}

% ══ RESUMEN GENERAL ══════════════════════════════════════════════════════════
\section*{Resumen general}
El estudiante demuestra un buen entendimiento de los principios de la mecánica de fluidos aplicados a bioanálisis. La mayoría de los problemas fueron resueltos correctamente, mostrando un buen manejo de fórmulas, cálculos y unidades. Sin embargo, se identificó un error conceptual significativo en la pregunta 4, donde no se consideró una de las variables clave del problema.

% ══ FORTALEZAS Y ÁREAS DE MEJORA ════════════════════════════════════════════
\vspace{4pt}
\begin{minipage}[t]{0.48\linewidth}
  \begin{tcolorbox}[cajaFortaleza, title={Fortalezas}]
\begin{itemize}[leftmargin=1.5em, itemsep=2pt]
  \item Excelente manejo de las ecuaciones de continuidad y Bernoulli, aplicándolas correctamente en un contexto de hemodinámica.
  \item Correcta aplicación del Principio de Arquímedes para determinar la densidad de un tejido.
  \item Precisión en los cálculos y consistencia en el uso de unidades en la mayoría de los problemas.
  \item Buena organización y claridad en la presentación de los datos y el desarrollo de las soluciones.
\end{itemize}
  \end{tcolorbox}
\end{minipage}
\hfill
\begin{minipage}[t]{0.48\linewidth}
  \begin{tcolorbox}[cajaMejora, title={Areas de mejora}]
\begin{itemize}[leftmargin=1.5em, itemsep=2pt]
  \item Análisis completo de los factores que influyen en las leyes físicas, especialmente en problemas de análisis proporcional donde múltiples variables cambian simultáneamente.
  \item Revisión exhaustiva de la interpretación de todas las condiciones dadas en el enunciado de los problemas para asegurar que ninguna variable sea omitida.
\end{itemize}
  \end{tcolorbox}
\end{minipage}

% ══ OBSERVACIONES POR PREGUNTA ═══════════════════════════════════════════════
\section*{Observaciones por pregunta}

\begin{longtable}{>{\bfseries}p{2.8cm} p{1.6cm} p{7.0cm} p{2.8cm}}
  \toprule
  \textbf{Pregunta} & \textbf{Puntaje} & \textbf{Evaluación} & \textbf{Errores detectados} \\
  \midrule
  \endhead
  \bottomrule
  \endfoot
    \textbf{Pregunta 1: Presión Hidrostática y Venosa} & 1.8/2.0 & El estudiante interpretó correctamente que se debía calcular la altura 'h' necesaria para generar la presión hidrostática requerida para vencer la presión venosa. El cálculo es correcto y las unidades son adecuadas. La notación inicial "P\_hidrostática = P\_manometric" podría ser más clara, pero el procedimiento subsiguiente es impecable para determinar la altura. & \textit{Ambigüedad menor en la notación inicial de la igualdad de presiones, aunque el cálculo final es correcto para la altura.} \\
    \midrule
    \textbf{Pregunta 2: Principio de Arquímedes} & 2.0/2.0 & La resolución es impecable. El estudiante calculó correctamente la masa del objeto, la fuerza de empuje, el volumen desplazado y, finalmente, la densidad del tejido óseo. Todos los pasos son lógicos, los cálculos precisos y las unidades consistentes. & \textit{---} \\
    \midrule
    \textbf{Pregunta 3: Hemodinámica (Continuidad y Bernoulli)} & 2.0/2.0 & La aplicación de la ecuación de continuidad para encontrar la velocidad en la estenosis es correcta. Posteriormente, la aplicación de la ecuación de Bernoulli para calcular la presión en la zona estrecha es precisa, incluyendo la consideración de la diferencia de velocidades. Los cálculos son exactos. & \textit{---} \\
    \midrule
    \textbf{Pregunta 4: Ley de Poiseuille (Análisis Proporcional)} & 0.5/2.0 & El estudiante solo consideró el efecto de la reducción del radio en el flujo (r\_f = 0.80r\_i, lo que lleva a un factor de (0.80)$^{4}$ = 0.4096). Sin embargo, omitió por completo el efecto del aumento de la viscosidad ($\eta$\_f = 1.80$\eta$\_i), que es una parte crucial del problema. La interpretación final del 40.96\% como el porcentaje de flujo es incorrecta, ya que solo representa el factor de cambio debido al radio. & \textit{Omisión del efecto del cambio en la viscosidad en la Ley de Poiseuille.; Interpretación incorrecta del resultado parcial (solo por cambio de radio) como el porcentaje final de flujo.} \\
    \midrule
    \textbf{Pregunta 5: Flujo Viscoso y Resistencia} & 2.0/2.0 & La resolución es completamente correcta. El estudiante identificó la fórmula adecuada (Ley de Poiseuille despejada para $\Delta$P), realizó las conversiones de unidades necesarias (diámetro a radio) y ejecutó los cálculos con precisión. El resultado final es correcto y con unidades apropiadas. & \textit{---} \\
    \midrule
\end{longtable}

% ══ ERRORES DETECTADOS ═══════════════════════════════════════════════════════
\section*{Análisis de errores}

\subsection*{Errores conceptuales}
\begin{itemize}[leftmargin=1.5em, itemsep=2pt]
  \item En la Pregunta 4, el estudiante no consideró todos los factores que afectan el flujo según la Ley de Poiseuille (específicamente, el cambio en la viscosidad), lo que llevó a un resultado incompleto y erróneo en el análisis proporcional.
\end{itemize}

\subsection*{Errores procedimentales}
\begin{itemize}[leftmargin=1.5em, itemsep=2pt]
  \item No se detectaron errores procedimentales significativos en los cálculos donde el concepto fue aplicado correctamente. El error en la Pregunta 4 fue más conceptual (omisión de una variable) que de cálculo.
\end{itemize}

% ══ RECOMENDACIONES AL ESTUDIANTE ════════════════════════════════════════════
\section*{Retroalimentación para el estudiante}
\begin{tcolorbox}[cajaRecomendacion, title={Dirigido a: rosmeilys sucre 33386357}]
Estimado/a estudiante, tu desempeño en este examen es bueno, demostrando un sólido conocimiento en la mayoría de los temas de mecánica de fluidos. Destacas en la aplicación de los principios de Arquímedes, Continuidad y Bernoulli, así como en la Ley de Poiseuille cuando todos los datos son explícitos. Sin embargo, es crucial que prestes más atención a la lectura completa de los enunciados, especialmente en problemas de análisis proporcional. En la Pregunta 4, la omisión del efecto de la viscosidad en el flujo fue un error importante. Asegúrate de identificar todas las variables que cambian y cómo afectan la magnitud física que se te pide calcular. Practicar más problemas de análisis de relaciones y proporciones te ayudará a fortalecer esta área. ¡Sigue así!
\end{tcolorbox}

% ══ NOTA PARA EL PROFESOR ════════════════════════════════════════════════════
\section*{Nota para el profesor}
\begin{tcolorbox}[
  enhanced, breakable,
  colback=yellow!8, colframe=yellow!60!black,
  fonttitle=\bfseries, coltitle=black,
  title={Observaciones para revisión manual},
  top=4pt, bottom=4pt, left=8pt, right=8pt,
]
El estudiante ha demostrado un buen nivel general en la resolución de problemas de mecánica de fluidos. El principal punto débil se encuentra en la Pregunta 4, donde omitió el efecto de la viscosidad en el análisis proporcional de la Ley de Poiseuille. Esto sugiere una necesidad de reforzar la comprensión de cómo múltiples variables interactúan en una fórmula y cómo realizar análisis de sensibilidad o proporcionales completos. El resto de las preguntas fueron resueltas con alta precisión.
\end{tcolorbox}

% ══ PIE DE INFORME ════════════════════════════════════════════════════════════
\vspace{12pt}
\begin{center}
  \footnotesize\color{gray}
  Informe generado automáticamente por el sistema de evaluación con IA (gemini-2.5-flash) y soporte LaTex. \\
  Este documento es una evaluación \textbf{preliminar} para ser revisado por el profesor antes de ser entregado al 
  estudiante. \\
  Generado el 2026-06-26 \textbullet\ BIOANALISIS Sistema Académico de Evaluación.
  \end{center}

\end{document}
